\documentclass{article}
\usepackage{amsmath}
\newcommand{\msun}{M_\odot}
\title{A Sample LaTeX Paper for the Phase 5 Pipeline}
\author{Jane Doe\thanks{E-mail: jane@example.org} \and John Roe}
\begin{document}
\maketitle
\begin{abstract}
We describe a minimal LaTeX source used to exercise the arXiv ingest pipeline.
\end{abstract}

\section{Introduction}\label{sec:intro}
As shown by \citet{KeyA} and \citep[see][p.~5]{KeyB}, the magnitude matters; see
also \citep{KeyA, KeyB}. The mass is $m = \msun$ and the cut was $|Z|<150\ $pc with
$\alpha\textsubscript{c}$ near $11\arcsec$. The energy follows Eq.~\eqref{eq:e},
discussed further in Sect.~\ref{sec:methods}, Fig.~\ref{fig:x}, and Table~\ref{tab:t}.
A note here.\footnote{This footnote must be dropped.} See \url{https://example.org}.

\begin{equation}\label{eq:e}
E = m c^2
\end{equation}
\begin{equation*}
x = 1
\end{equation*}
\begin{align}
a &= b \\
c &= d \label{eq:al}
\end{align}
\[ y = 2 \]

\section{Methods}\label{sec:methods}
We refer back to Sect.~\ref{sec:intro}.

\begin{figure}
\centering
\includegraphics[width=\linewidth]{nonexistent_figure.pdf}
\caption{A sample figure caption citing \citet{KeyA}.}
\label{fig:x}
\end{figure}

\begin{table}
\caption{A sample table caption.}
\label{tab:t}
\begin{tabular}{ll}
Header A & Header B \\
1 & 2 \\
\end{tabular}
\end{table}

\begin{thebibliography}{9}
\bibitem[Smith et al.(2020)]{KeyA} Smith, J., Doe, J., \& Roe, J. 2020, \mnras, 600, 1
  \doi{10.1051/0004-6361/200000001}
\bibitem[Jones \& Roe(2019)]{KeyB} Jones, A., \& Roe, B. 2019, \aap, 590, 12
\end{thebibliography}
\end{document}
